% LTspice Transient Waveform Plot for URN Paper
% Generated from actual LTspice simulation results
% Font: Times New Roman, 10pt
% Tick direction: inward (box on)

\documentclass[10pt]{article}
\usepackage[T1]{fontenc}
\usepackage{mathptmx}  % Times New Roman
\usepackage{pgfplots}
\pgfplotsset{compat=1.18}
\usepgfplotslibrary{units}

\begin{document}

% Define common style
\pgfplotsset{
    every axis/.append style={
        font=\fontsize{10}{12}\selectfont,
        tick align=inside,  % Tick marks inside (inward)
        major tick length=3pt,
        minor tick length=1.5pt,
        line width=0.8pt,
        axis line style={line width=0.8pt},
        tick style={line width=0.6pt},
        grid style={line width=0.3pt, gray!30},
    }
}

%% Figure (a): Input Voltage and Current
\begin{figure}[htbp]
\centering
\begin{tikzpicture}
\begin{axis}[
    width=0.85\textwidth,
    height=0.45\textwidth,
    xlabel={Time (ms)},
    ylabel={Input Voltage $V_\mathrm{in}$ (V)},
    xmin=0, xmax=20,
    ymin=-0.1, ymax=1.2,
    xtick={0,5,10,15,20},
    ytick={0,0.5,1.0},
    grid=major,
    axis y line*=left,
    legend style={at={(0.02,0.98)}, anchor=north west, font=\footnotesize},
]
\addplot[blue, thick, no markers] table[
    x=time_ms,
    y=V_in,
    col sep=comma,
    comment chars={\#},
] {ltspice_waveform_data.csv};
\addlegendentry{$V_\mathrm{in}$}
\end{axis}

\begin{axis}[
    width=0.85\textwidth,
    height=0.45\textwidth,
    xmin=0, xmax=20,
    ymin=-400, ymax=400,
    ylabel={Input Current $I_\mathrm{in}$ (mA)},
    axis y line*=right,
    axis x line=none,
    ytick={-300,-150,0,150,300},
    legend style={at={(0.98,0.98)}, anchor=north east, font=\footnotesize},
    tick align=inside,
]
\addplot[green!60!black, thick, no markers] table[
    x=time_ms,
    y=I_in_mA,
    col sep=comma,
    comment chars={\#},
] {ltspice_waveform_data.csv};
\addlegendentry{$I_\mathrm{in}$}
\end{axis}
\end{tikzpicture}
\caption{Transient response of URN equivalent circuit: input voltage pulse and resulting current.
\textbf{Data source: Actual LTspice simulation} (not analytical approximation).}
\label{fig:ltspice_transient_input}
\end{figure}

%% Figure (b): Output Voltage (scaled for visibility)
\begin{figure}[htbp]
\centering
\begin{tikzpicture}
\begin{axis}[
    width=0.85\textwidth,
    height=0.45\textwidth,
    xlabel={Time (ms)},
    ylabel={Output Voltage $V_\mathrm{out}$ (mV)},
    xmin=0, xmax=20,
    ymin=0, ymax=0.4,
    xtick={0,5,10,15,20},
    ytick={0,0.1,0.2,0.3,0.4},
    scaled y ticks=false,
    yticklabel style={/pgf/number format/fixed, /pgf/number format/precision=1},
    grid=major,
    legend style={at={(0.98,0.02)}, anchor=south east, font=\footnotesize},
]
% Note: V_out in CSV is in V, multiply by 1000 for mV display
\addplot[red, thick, no markers] table[
    x=time_ms,
    y expr=\thisrow{V_out}*1000,
    col sep=comma,
    comment chars={\#},
] {ltspice_waveform_data.csv};
\addlegendentry{$V_\mathrm{out}$ (LTspice)}

% Add annotation for no ringing
\node[fill=green!20, draw=green!60!black, rounded corners, font=\footnotesize]
    at (axis cs:16,0.32) {No spurious ringing};
\end{axis}
\end{tikzpicture}
\caption{Output voltage of URN RLC ladder circuit showing smooth charging characteristic.
The monotonic response confirms the passivity-by-construction property of URN-synthesized circuits.
\textbf{Data source: Actual LTspice simulation via PyLTSpice API.}}
\label{fig:ltspice_transient_output}
\end{figure}

%% Combined single figure for paper
\begin{figure}[htbp]
\centering
\begin{tikzpicture}
% Input voltage subplot
\begin{axis}[
    name=plot1,
    width=0.48\textwidth,
    height=0.35\textwidth,
    xlabel={Time (ms)},
    ylabel={Voltage (V)},
    xmin=0, xmax=20,
    ymin=-0.1, ymax=1.2,
    xtick={0,5,10,15,20},
    ytick={0,0.5,1.0},
    grid=major,
    title={\footnotesize (a) Input voltage pulse},
    title style={at={(0.5,1.0)}, anchor=south},
]
\addplot[blue, thick, no markers] table[
    x=time_ms,
    y=V_in,
    col sep=comma,
    comment chars={\#},
] {ltspice_waveform_data.csv};
\end{axis}

% Input current subplot
\begin{axis}[
    name=plot2,
    at={(plot1.south east)},
    anchor=south west,
    xshift=0.8cm,
    width=0.48\textwidth,
    height=0.35\textwidth,
    xlabel={Time (ms)},
    ylabel={Current (mA)},
    xmin=0, xmax=20,
    ymin=-400, ymax=400,
    xtick={0,5,10,15,20},
    ytick={-300,0,300},
    grid=major,
    title={\footnotesize (b) Input current response},
    title style={at={(0.5,1.0)}, anchor=south},
]
\addplot[green!60!black, thick, no markers] table[
    x=time_ms,
    y=I_in_mA,
    col sep=comma,
    comment chars={\#},
] {ltspice_waveform_data.csv};
\end{axis}
\end{tikzpicture}
\caption{LTspice transient simulation results of URN-synthesized RLC ladder circuit.
(a) Input voltage pulse (0--1~V, 5~ms period).
(b) Input current showing RC charging/discharging behavior without spurious oscillations.
\textbf{Note: These are actual LTspice simulation results, not analytical Python approximations.}}
\label{fig:ltspice_combined}
\end{figure}

\end{document}
